\section{Direct mechanical identification and optional noisy observation}\label{sec:observation}
\subsection{The collision bit already carries regular information}
Either the prepared-tube probability $p(R)$ on $I_0$ or the equilibrium probability $p_{\mathrm{eq}}(R)$ on $I$ defines a deterministic binary readout $C_R$. In each case $p$ is smooth, is bounded away from zero and one on its compact parameter interval, and $p'>0$.

\begin{theorem}[Identification without Gaussian noise]\label{thm:binary-info}
The experiment $P_R(C=1)=p(R)$ is globally identifiable on the stated interval. Its score and Fisher information are
\[
 \ell_R(C)=\frac{p'(R)(C-p(R))}{p(R)(1-p(R))},\qquad
 \mathcal I_C(R)=\frac{p'(R)^2}{p(R)(1-p(R))}>0.
\]
At every interior parameter it is differentiable in quadratic mean, with squared remainder $O(h^4)$ uniformly on compact subintervals when $p''$ is bounded.
\end{theorem}
\begin{proof}
Strict monotonicity of $p$ gives identification. Differentiate the two positive probabilities: the scores are $p'/p$ at $C=1$ and $-p'/(1-p)$ at $C=0$. Their mean is zero and their squared mean gives the displayed information. The square-root probability vector
$(\sqrt{1-p(R)},\sqrt{p(R)})$ has a bounded second derivative on the compact interval. Taylor's integral remainder has Euclidean norm $O(h^2)$; squaring gives $O(h^4)$. Its first derivative is exactly half the score times the square-root vector, proving quadratic-mean differentiability. The only randomness is the explicitly stated initial preparation; no measurement-noise distribution is needed.
\end{proof}

This is a finite experiment, not a theorem about infinitely many dependent collisions. Independent repetition requires independent re-preparation. A repeated measurement along one retained trajectory has a different likelihood.

\subsection{Adding a channel loses, rather than creates, this information}
Let the apparatus produce a parameter-independent observation channel $g_c(y)$, $c=0,1$, each a normalized density relative to $\lambda$. The observed density is
\[
 r_R(y)=(1-p(R))g_0(y)+p(R)g_1(y).
\]

\begin{theorem}[Observed information and quadratic-mean response]\label{thm:channel-info}
If $g_0\ne g_1$ in $L^1$, this experiment remains identifiable. Its Fisher information satisfies
\begin{align}
 \mathcal I_Y(R)&=p'(R)^2\int\frac{(g_1-g_0)^2}{r_R}\dd\lambda,\\
 p'(R)^2\norm{g_1-g_0}_1^2&\le\mathcal I_Y(R)
 \le\frac{p'(R)^2}{p(R)(1-p(R))}.
 \label{eq:channel-info}
\end{align}
It is differentiable in quadratic mean with squared remainder $O(h^4)$ on compact parameter intervals. Gaussian translated densities are one optional instance, not an assumption of the mechanical theorem.
\end{theorem}
\begin{proof}
Equality of two observed densities implies $(p(R)-p(\widetilde R))(g_1-g_0)=0$, hence $R=\widetilde R$. The derivative is $p'(g_1-g_0)$, giving the information formula. Cauchy--Schwarz applied with factors $|g_1-g_0|/\sqrt{r_R}$ and $\sqrt{r_R}$ proves the lower bound. The observed score is $E[\ell_R(C)\mid Y]$, so conditional Jensen and Theorem~\ref{thm:binary-info} prove the upper bound.

Let $c_*\le\min(p,1-p)$ on the interval. Then $r_R\ge c_*(g_0+g_1)$, and $|r_R'|+|r_R''|\le C(g_0+g_1)$. The second derivative of the square root,
$r_R''/(2\sqrt{r_R})-(r_R')^2/(4r_R^{3/2})$, is bounded by $C\sqrt{g_0+g_1}$, which lies in $L^2(\lambda)$. Banach-valued Taylor expansion proves the claimed squared remainder. Possible points where both densities vanish contribute nothing and are assigned zero square-root derivatives.
\end{proof}

\subsection{Smooth observation of a singular mechanical law}
Let $X_a$ be a mechanical output with a finite branch decomposition or a completed assembly. Suppose the observation density $g(y\mid x)$ satisfies on every relevant compact output set
\[
 \int\sup_x\sum_{|\alpha|\le r}|\partial_x^\alpha g(y\mid x)|\lambda(dy)<\infty.
\]

\begin{proposition}[From currents to strong observation response]\label{prop:smoothing}
If the mechanical branch transports, weights, and output derivatives through order $r$ are uniformly bounded and the complete assembly has a common derivative envelope, then $r_a(y)=E[g(y\mid X_a)]$ is $C^r$ in $L^1(\lambda)$. The same holds for a numerator $N_a(y)=E[F_ag(y\mid X_a)]$ when the test family has the matching bounds. Mechanical interface jumps are included in these derivatives.
\end{proposition}
\begin{proof}
Pull each regular branch to a fixed domain, or use the fixed-domain tube assembly of Section~\ref{sec:lorentz-global}. Each derivative is a finite combination of channel derivatives, output derivatives, and weight derivatives. The displayed channel envelope and the mechanical bounds dominate the integrals jointly in output and preparation coordinates. The Banach-valued fundamental theorem of calculus proves $L^1$ differentiability. Sum the completed assemblies using their derivative envelope. In physical coordinates the same differentiation produces the interface terms of Theorem~\ref{thm:shape}; they are not lost by observation smoothing.
\end{proof}

\subsection{Evidence-weighted prediction and null sets}
For $|F_a|\le B$, choose $u_a=N_a/r_a$ where $r_a>0$ and a version with $|u_a|\le B$ elsewhere. For two experiments,
\begin{equation}
 \int\widetilde r|\widetilde u-u|\dd\lambda
 \le\norm{\widetilde N-N}_1+B\norm{\widetilde r-r}_1.
 \label{eq:evidence}
\end{equation}
Indeed $N=ru$ and $\widetilde N=\widetilde r\widetilde u$ hold everywhere, so subtract and apply the triangle inequality. If pointwise differentiable versions are supplied, then $ru'=N'-ur'$ and
$\int r|u'|\le\norm{N'}_1+B\norm{r'}_1$.
An $L^1$ derivative alone does not assert pointwise differentiability; the smooth channel envelope supplies it where used. A bound at each time is not a bound for the expected path supremum or a theorem on changing filtrations.
\par
